% ============================================================
% 第3章补充：经典例题与自学元素
% 插入位置：第3章现有miniquiz之后（\end{miniquiz} 之后）
% ============================================================

\section*{习题精选与详解}
\addcontentsline{toc}{section}{习题精选}

\begin{example}{托卡马克中不同粒子的回旋参数}{gyro-params}
在磁约束聚变装置中，磁场 $B=5\,\mathrm{T}$，等离子体温度 $T=10\,\mathrm{keV}$。分别计算电子、质子（$\mathrm{H}^+$）和氘核（$\mathrm{D}^+$）的回旋频率、回旋半径和回旋周期。质量取 $m_e=9.11\times 10^{-31}\,\mathrm{kg}$，$m_p=1.673\times 10^{-27}\,\mathrm{kg}$，$m_D\approx 3.34\times 10^{-27}\,\mathrm{kg}$。
\tcblower
\textbf{解答：}
$k_{\!B}T = 10^4\,\mathrm{eV} = 1.602\times 10^{-15}\,\mathrm{J}$。热速度 $v_{th}=\sqrt{2k_{\!B}T/m}$。

\textbf{电子：}
\[
v_{th,e} = \sqrt{\frac{2\times 1.602\times 10^{-15}}{9.11\times 10^{-31}}} = \sqrt{3.51\times 10^{15}} \approx 5.93\times 10^{7}\,\mathrm{m/s}
\]
\[
\omega_{ce} = \frac{eB}{m_e} = \frac{1.602\times 10^{-19}\times 5}{9.11\times 10^{-31}} \approx 8.79\times 10^{11}\,\mathrm{rad/s}
\]
\[
f_{ce} = \frac{\omega_{ce}}{2\pi} \approx 140\,\mathrm{GHz},\qquad r_{ce} = \frac{v_{th,e}}{\omega_{ce}} \approx \frac{5.93\times 10^{7}}{8.79\times 10^{11}} \approx 6.7\times 10^{-5}\,\mathrm{m} = 67\,\mu\mathrm{m}
\]

\textbf{质子：}
\[
v_{th,p} = \sqrt{\frac{2\times 1.602\times 10^{-15}}{1.673\times 10^{-27}}} \approx 1.38\times 10^{6}\,\mathrm{m/s}
\]
\[
\omega_{cp} = \frac{eB}{m_p} = \frac{1.602\times 10^{-19}\times 5}{1.673\times 10^{-27}} \approx 4.79\times 10^{8}\,\mathrm{rad/s}
\]
\[
f_{cp} \approx 76.2\,\mathrm{MHz},\qquad r_{cp} \approx \frac{1.38\times 10^{6}}{4.79\times 10^{8}} \approx 2.88\times 10^{-3}\,\mathrm{m} = 2.88\,\mathrm{mm}
\]

\textbf{氘核：}
\[
v_{th,D} = \sqrt{\frac{2\times 1.602\times 10^{-15}}{3.34\times 10^{-27}}} \approx 9.79\times 10^{5}\,\mathrm{m/s}
\]
\[
\omega_{cD} = \frac{eB}{m_D} \approx 2.40\times 10^{8}\,\mathrm{rad/s},\quad f_{cD}\approx 38.2\,\mathrm{MHz},\quad r_{cD}\approx 4.08\,\mathrm{mm}
\]

\textbf{对比：} 电子回旋频率比离子高约 $10^3$ 倍，回旋半径则小约 $10^2$ 倍。这就是为什么在磁约束中，电子通常被很好约束，而离子轨道半径较大，需更大的装置尺寸。
\end{example}

\begin{example}{$E\times B$ 漂移速度计算}{exb-drift-calc}
在均匀电磁场中，$\vect{B}=5\,\mathrm{T}\,\unitvec{z}$，$\vect{E}=10^3\,\mathrm{V/m}\,\unitvec{y}$。计算正离子和电子的 $E\times B$ 漂移速度，并判断其是否与粒子质量、电荷量或速度无关。
\tcblower
\textbf{解答：}
$E\times B$ 漂移速度公式为
\[
\vect{v}_E = \frac{\vect{E}\times\vect{B}}{B^2}
\]

先计算 $\vect{E}\times\vect{B}$：
\[
\vect{E}\times\vect{B} = (E\unitvec{y})\times(B\unitvec{z}) = EB\,(\unitvec{y}\times\unitvec{z}) = EB\,\unitvec{x}
\]

代入数值：
\[
\vect{v}_E = \frac{10^3\times 5}{5^2}\,\unitvec{x} = \frac{5\times 10^3}{25}\,\unitvec{x} = 200\,\mathrm{m/s}\,\unitvec{x}
\]

\textbf{讨论：} 结果 $200\,\mathrm{m/s}$ 与粒子质量、电荷量、速度大小均无关。这意味着等离子体中所有粒子（正离子、电子、不同质量离子）都以相同速度向 $+x$ 方向漂移。这是 $E\times B$ 漂移的重要特征，也是磁约束装置中粒子输运的基础之一。
\end{example}

\begin{example}{磁镜损失锥与粒子反射判断}{mirror-loss-cone}
某磁镜装置中，轴上最小磁场 $B_{\min}=1\,\mathrm{T}$，磁镜喉部最大磁场 $B_{\max}=3\,\mathrm{T}$。
\begin{enumerate}[label=(\arabic*)]
\item 求损失锥半角 $\theta_m$；
\item 一粒子在 $B_{\min}$ 处总速度 $v=2\times 10^6\,\mathrm{m/s}$，平行速度 $v_{\parallel}=1.8\times 10^6\,\mathrm{m/s}$，判断该粒子能否被约束。
\end{enumerate}
\tcblower
\textbf{解答：}

\textbf{(1) 损失锥半角：}
由磁镜比 $R_m = B_{\max}/B_{\min}=3$，损失锥条件为
\[
\sin^2\theta_m = \frac{1}{R_m} = \frac{1}{3}
\]
\[
\theta_m = \arcsin\!\left(\frac{1}{\sqrt{3}}\right) \approx 35.3^{\circ}
\]

\textbf{(2) 粒子判断：}
先求垂直速度：
\[
v_{\perp} = \sqrt{v^2 - v_{\parallel}^2} = \sqrt{(2\times 10^6)^2 - (1.8\times 10^6)^2} = \sqrt{0.76\times 10^{12}} \approx 8.72\times 10^{5}\,\mathrm{m/s}
\]

计算投掷角的正弦平方：
\[
\sin^2\theta = \left(\frac{v_{\perp}}{v}\right)^2 = \left(\frac{8.72\times 10^5}{2\times 10^6}\right)^2 \approx (0.436)^2 \approx 0.19
\]

与损失锥条件比较：
\[
\sin^2\theta = 0.19 < \frac{1}{3} \approx 0.333
\]

\textbf{结论：} 该粒子的 $\sin^2\theta$ 小于 $1/R_m$，位于损失锥内，因此它将穿过磁镜喉部逃逸，无法被约束。要约束该粒子，需要增大磁镜比（提高 $B_{\max}$）或注入具有更大垂直速度分量的粒子。
\end{example}

\begin{example}{托卡马克中的梯度漂移与曲率漂移}{tokamak-drifts}
在一大环半径 $R=2\,\mathrm{m}$ 的托卡马克中，磁场 $B=5\,\mathrm{T}$。取电子温度与离子温度均为 $T=10\,\mathrm{keV}$，且速度分布各向同性。分别估算电子和氘核的 $E\times B$ 漂移、梯度漂移与曲率漂移的合速度大小与方向。
\tcblower
\textbf{解答：}
在圆环磁场中，$|\nabla B|\approx B/R$，且梯度漂移与曲率漂移方向相同（垂直环面）。两者可合并为
\[
\vect{v}_{\nabla B+R} = \frac{m}{qB^2R}\left(v_{\parallel}^2 + \frac{v_{\perp}^2}{2}\right)\hat{\vect{R}}\times\vect{B}
\]

对各向同性分布，$\langle v_{\parallel}^2\rangle = v^2/3$，$\langle v_{\perp}^2\rangle = 2v^2/3$，于是
\[
\langle v_{\parallel}^2 + v_{\perp}^2/2\rangle = v^2/3 + v^2/3 = \frac{2v^2}{3}
\]

故漂移速度大小为
\[
v_{\text{drift}} = \frac{2mv^2}{3qB^2R} = \frac{4k_{\!B}T}{3qB^2R}
\]
最后一步利用了 $mv^2 = 3k_{\!B}T$（三维各向同性）。

\textbf{电子漂移：}
\[
v_{\text{drift},e} = \frac{4\times 1.602\times 10^{-15}}{3\times (-1.602\times 10^{-19})\times 25\times 2} = \frac{6.408\times 10^{-15}}{-2.403\times 10^{-17}} \approx -267\,\mathrm{m/s}
\]
方向垂直环面向\textbf{下}（负号表示与离子反向）。

\textbf{氘核漂移：}
\[
v_{\text{drift},D} = \frac{4\times 1.602\times 10^{-15}}{3\times (1.602\times 10^{-19})\times 25\times 2} \approx +267\,\mathrm{m/s}
\]
方向垂直环面向\textbf{上}。

\textbf{物理后果：} 电子与离子漂移方向相反，在托卡马克中形成垂直环面的电荷分离，产生电场。该电场又驱动 $E\times B$ 漂移，形成环向输运（即经典扩散与新经典输运的根源之一）。漂移大小仅约 $270\,\mathrm{m/s}$，远低于热速度，但长期累积造成显著的粒子与能量损失。
\end{example}

\begin{figure}[htbp]
\centering
\begin{tikzpicture}[scale=1.0]
% 左图：正离子
\begin{scope}[xshift=-4.5cm]
\draw[->,thick] (-2,0) -- (2.5,0) node[right] {$x$};
\draw[->,thick] (0,-1.5) -- (0,2) node[above] {$y$};

% 磁场出纸面
\foreach \x in {-1.5,-0.75,0,0.75,1.5} {
  \foreach \y in {-1,-0.5,0,0.5,1} {
    \node at (\x,\y) {\tiny $\odot$};
  }
}
\node[right] at (1.5,1.5) {$\vect{B}\parallel +\unitvec{z}$};

% 电场向上
\draw[->,red,thick] (0,1.5) -- (0,2.3) node[right] {$\vect{E}$};

% 正离子轨迹：逆时针回旋 + 向右漂移
\draw[blue,thick,domain=0:720,samples=200,smooth] plot({0.7*cos(\x)+0.08*\x/180},{0.7*sin(\x)});
\fill[blue] (0.7,0) circle (2pt);
\node[blue,align=center] at (0.8,-1.2) {正离子\\$r_{ci}\sim 1\,\mathrm{mm}$};
\draw[->,blue,thick] (0.7,0) -- (0.75,0.2);

% 漂移方向
\draw[->,blue!60!black,very thick] (1.5,-1.5) -- (2.3,-1.5) node[right] {$\vect{v}_E$};
\end{scope}

% 右图：电子
\begin{scope}[xshift=3cm]
\draw[->,thick] (-2,0) -- (2.5,0) node[right] {$x$};
\draw[->,thick] (0,-1.5) -- (0,2) node[above] {$y$};

% 磁场出纸面
\foreach \x in {-1.5,-0.75,0,0.75,1.5} {
  \foreach \y in {-1,-0.5,0,0.5,1} {
    \node at (\x,\y) {\tiny $\odot$};
  }
}
\node[right] at (1.5,1.5) {$\vect{B}\parallel +\unitvec{z}$};

% 电场向上
\draw[->,red,thick] (0,1.5) -- (0,2.3) node[right] {$\vect{E}$};

% 电子轨迹：顺时针回旋 + 向右漂移
\draw[orange,thick,domain=0:720,samples=200,smooth] plot({0.2*cos(-\x)+0.08*\x/180},{0.2*sin(-\x)});
\fill[orange] (0.2,0) circle (1.5pt);
\node[orange,align=center] at (0.3,-0.6) {电子\\$r_{ce}\sim 70\,\mu\mathrm{m}$};
\draw[->,orange,thick] (0.2,0) -- (0.22,-0.08);

% 漂移方向
\draw[->,orange!60!black,very thick] (1.5,-1.5) -- (2.3,-1.5) node[right] {$\vect{v}_E$};
\end{scope}
\end{tikzpicture}
\caption{$E\times B$ 漂移的摆线轨迹。磁场垂直纸面向外，电场沿 $+y$ 方向。正离子逆时针回旋，电子顺时针回旋，但两者漂移方向相同（均沿 $+x$）。}
\end{figure}

\begin{figure}[htbp]
\centering
\begin{tikzpicture}[scale=1.1]
% 磁镜轮廓（沙漏形）
\draw[thick,blue] (-2,3) .. controls (-1,2) and (-0.7,1) .. (-0.5,0.5);
\draw[thick,blue] (-2,-3) .. controls (-1,-2) and (-0.7,-1) .. (-0.5,-0.5);
\draw[thick,blue] (2,3) .. controls (1,2) and (0.7,1) .. (0.5,0.5);
\draw[thick,blue] (2,-3) .. controls (1,-2) and (0.7,-1) .. (0.5,-0.5);
\draw[thick,blue] (-0.5,0.5) -- (-0.5,-0.5);
\draw[thick,blue] (0.5,0.5) -- (0.5,-0.5);

% 磁场标注
\node[blue,right] at (2,2.5) {$B_{\max}$};
\node[blue] at (0,0) {$B_{\min}$};
\draw[->,blue] (1.8,2.8) -- (1.2,2.2);
\draw[->,blue] (-1.8,2.8) -- (-1.2,2.2);
\draw[->,blue] (1.8,-2.8) -- (1.2,-2.2);
\draw[->,blue] (-1.8,-2.8) -- (-1.2,-2.2);

% 损失锥边界
\draw[red,thick] (-1.5,2.5) -- (0,0) -- (1.5,2.5);
\draw[red,thick] (-1.5,-2.5) -- (0,0) -- (1.5,-2.5);
\draw[red,->] (0,0.4) arc (90:55:0.4);
\draw[red,->] (0,-0.4) arc (-90:-55:0.4);
\node[red] at (0.25,1.5) {$\theta_m$};
\node[red] at (0.25,-1.5) {$\theta_m$};

% 被反射粒子（大角度）
\draw[green!60!black,thick,->] (0,-2.2) -- (0.2,-1.2) -- (0,-0.3);
\draw[green!60!black,thick,->] (0,-0.3) -- (-0.2,0.8) -- (0,1.8);
\node[green!60!black,right] at (0.4,-1.5) {被反射};
\node[green!60!black] at (-0.6,1.0) {\small $\theta>\theta_m$};

% 逃逸粒子（小角度）
\draw[orange,thick,->] (0,-2.2) -- (0.1,-1.2) -- (0.2,-0.3);
\draw[orange,thick,->] (0.2,-0.3) -- (0.5,0.8) -- (1.0,2.2);
\node[orange,right] at (1.2,1.5) {逃逸};
\node[orange] at (0.8,0.2) {\small $\theta<\theta_m$};
\end{tikzpicture}
\caption{磁镜损失锥示意图。在磁镜喉部 $B_{\max}$ 处，若粒子投掷角 $\theta<\theta_m$，则 $v_{\perp}^2$ 不足以将平行动能完全转化为垂直动能，粒子将穿过磁镜逃逸。}
\end{figure}

% ch3 新例题：磁镜与损失锥
\begin{example}{磁镜约束与损失锥角估算}{ch3-mirror-est}
一磁镜装置中，中心磁场 $B_0 = 0.2\,\mathrm{T}$，镜点磁场 $B_m = 1.0\,\mathrm{T}$（磁镜比 $R_m = B_m/B_0 = 5$）。
（1）推导损失锥角 $\theta_{lc}$（以 $B_0$ 处计）的表达式并求其数值；
（2）若粒子在中心处垂直速度分量占总动能的 $90\%$（即 $\sin^2\theta = 0.9$），判断它是否被约束；
（3）说明磁镜约束为何无法完全约束各向同性等离子体。
\tcblower
\textbf{解答：}

\textbf{(1) 损失锥角：} 磁矩不变量 $\mu = m v_\perp^2/(2B)$ 要求 $v_\perp^2/B = \text{const}$，粒子向强场区运动时 $v_\perp$ 增大、$v_\parallel$ 减小。粒子被反射的条件是平行速度在到达镜点前降为零。结合能量守恒，粒子能到达镜点的条件为
\[
\frac{v_\perp^2}{v^2}\bigg|_{B_0} \le \frac{B_0}{B_m} = \frac{1}{R_m}
\]
不满足上式的粒子构成\textbf{损失锥}，锥角满足
\[
\sin^2\theta_{lc} = \frac{1}{R_m} \quad\Rightarrow\quad \theta_{lc} = \arcsin\!\left(\frac{1}{\sqrt{R_m}}\right) = \arcsin\!\left(\frac{1}{\sqrt{5}}\right) \approx 26.6^\circ
\]

\textbf{(2) 判断约束：} 该粒子 $\sin^2\theta = 0.9 > 1/R_m = 0.2$，$\theta > \theta_{lc}$，\textbf{在损失锥之外，被磁镜约束}。它在中心处 $v_\parallel^2/v^2 = 0.1$，到达镜点前平行速度降为零并反射。

\textbf{(3) 无法完全约束的原因：} 对各向同性分布（$\theta$ 均匀分布在立体角内），落入损失锥（$0\le\theta\le\theta_{lc}$）的粒子比例约为 $\sin^2\theta_{lc}/2 = 1/(2R_m) = 10\%$。每一轮反射都会损失锥内粒子，若无其他机制填充损失锥，等离子体将快速损失。磁镜对静电与湍流散射（使粒子扩散进损失锥）也很敏感。因此单纯磁镜无法稳定约束热等离子体，现代磁镜实验需辅以中性束注入、偶极场或 ECRH 填充等手段。
\end{example}

\begin{formula}{第3章核心公式速查}{ch3-formulas}
\begin{itemize}[leftmargin=*,itemsep=0pt]
\item 回旋频率：$\displaystyle\omega_c = \frac{|q|B}{m}$，$\displaystyle r_c = \frac{v_{\perp}}{\omega_c}$
\item $E\times B$ 漂移：$\displaystyle\vect{v}_E = \frac{\vect{E}\times\vect{B}}{B^2}$（与 $q,m$ 无关）
\item 梯度漂移：$\displaystyle\vect{v}_{\nabla B} = \frac{m v_{\perp}^2}{2qB^3}\,(\vect{B}\times\nabla B)$
\item 曲率漂移：$\displaystyle\vect{v}_R = \frac{m v_{\parallel}^2}{qB^2R^2}\,(\vect{R}\times\vect{B})$
\item 磁镜反射条件：$\displaystyle\sin^2\theta_{\min} > \frac{B_{\min}}{B_{\max}} = \frac{1}{R_m}$
\item 损失锥半角：$\displaystyle\theta_m = \arcsin\!\left(\frac{1}{\sqrt{R_m}}\right)$
\item 磁矩绝热不变量：$\displaystyle\mu = \frac{m v_{\perp}^2}{2B}$
\item 纵向绝热不变量：$\displaystyle J = \oint v_{\parallel}\,\diff s$
\end{itemize}
\end{formula}

\begin{checklist}{第3章自学检查清单}{ch3-checklist}
\begin{itemize}[leftmargin=*,label=$\square$]
\item 我能计算不同粒子在给定磁场中的回旋频率、半径和周期
\item 我掌握了 $E\times B$ 漂移公式，并理解其方向与粒子种类无关
\item 我能用梯度漂移和曲率漂移公式估算托卡马克中的粒子漂移
\item 我能计算磁镜损失锥角，并判断给定速度粒子是否被反射
\item 我理解了磁矩 $\mu$ 作为绝热不变量的条件：$\tau_{\text{bounce}}\ll\tau_{\text{slow}}$
\item 我知道为什么托卡马克中梯度漂移和曲率漂移会导致电荷分离
\end{itemize}
\end{checklist}

\endinput
